% CHƯƠNG 5: sáu thành phần cố định cho mỗi nội dung; hình/bảng chèn tự động.
\clearpage
\section[Kiểm chứng chuỗi không dừng và mô hình ARIMA]{KIỂM CHỨNG CHUỖI KHÔNG DỪNG VÀ MÔ HÌNH ARIMA}

Chương này kiểm chứng \muc{2.2} trên ba bộ dữ liệu không dừng là NI, AP và DP, trong đó AP được lấy log và chỉ dùng phần huấn luyện. DP trôi như một bước ngẫu nhiên, NI đổi mức quanh năm 1899, còn AP vừa có xu thế vừa có mùa vụ. Cách đọc giống \chg{4}, và các mô hình ARIMA khớp ở đây được dùng để dự báo ở \chg{6}.

\subsection{Chuỗi không dừng}

\noidung{Xu thế tất định và ngẫu nhiên}
\tp{Nội dung lý thuyết cần kiểm chứng} Kết quả \ndlk{2.2.1(a)} nói rằng trừ xu thế bằng hồi quy không làm hết tính dai dẳng của chuỗi $I(1)$, còn lấy sai phân thì đưa chuỗi $I(1)$ về dạng dừng.
\tp{Cách triển khai} Nội dung này dùng dữ liệu thực của NI, AP và DP. Với mỗi chuỗi, ta so chuỗi gốc, chuỗi đã trừ xu thế tuyến tính (bình phương tối thiểu) và sai phân bậc một, qua độ lệch chuẩn và $\hat\rho$ ở độ trễ 1 và 10.
\tp{Mã lệnh Python}
\begin{codeblock}
for ten, y in SERIES.items():                      # NI, AP (log, huan luyen), DP (log)
    n = len(y); t = np.arange(n); d1 = np.diff(y)
    khu = y - np.polyval(np.polyfit(t, y, 1), t)             # tru xu the tuyen tinh (OLS)
    sd_y, sd_d1 = y.std(ddof=1), d1.std(ddof=1)
    rho_y, rho_khu, rho_d1 = (A.sample_acf(v, 10) for v in (y, khu, d1))   # lay do tre 1 va 10
\end{codeblock}
\tp{Kết quả tính toán} Bảng~\ref{tab:tab_b12_01_levels_vs_differences} so độ lệch chuẩn và tự tương quan của chuỗi gốc, chuỗi trừ xu thế và sai phân, còn Hình~\ref{fig:fig_b12_01_levels_vs_differences} vẽ chuỗi cùng ACF của chuỗi gốc và của sai phân.
\input{\tabdir/tab_b12_01_levels_vs_differences.tex}
\figx{fig_b12_01_levels_vs_differences}{Chuỗi, ACF của chuỗi gốc và ACF của sai phân bậc một cho NI, AP và DP; AP và DP lấy log, AP chỉ gồm phần huấn luyện}
\tp{Phân tích và diễn giải}
\bodl{NI.} Chuỗi không có xu thế dài mà chỉ đổi mức quanh năm 1899. Sai phân gần như không giảm độ lệch chuẩn, chỉ từ $169.2$ xuống $168.1$, và đưa $\hat\rho(1)$ về $-0.40$, một giá trị âm đáng kể, dấu hiệu của sai phân thừa sẽ được xét ở \muc{5.3}. NI không thuộc hẳn loại nào trong hai loại xu thế.

\bodl{AP.} $\hat\rho(1)=0.95$ và $\hat\rho(10)=0.71$. Trừ xu thế kéo $\hat\rho(10)$ xuống $0.18$, nghĩa là phần lớn tính dai dẳng đến từ xu thế, và sai phân giảm độ lệch chuẩn từ $0.39$ xuống $0.11$. Nhưng ACF của sai phân vẫn có đỉnh lớn ở độ trễ 12 và 24, tức mùa vụ chưa được xử lý, như sẽ xét ở \muc{5.2}. Chỉ riêng ACF khó phân biệt hai loại xu thế ở đây.

\bodl{DP.} Đây là ví dụ rõ nhất của xu thế ngẫu nhiên: $\hat\rho(1)=0.998$ và $\hat\rho(10)=0.979$, và trừ xu thế gần như không thay đổi gì, chỉ còn $0.996$ và $0.960$. Sai phân đưa chuỗi về gần nhiễu trắng với $\hat\rho(1)=-0.0004$, và độ lệch chuẩn giảm từ $0.363$ xuống $0.0103$.

\tp{Nhận xét} Kết quả \textbf{phù hợp một phần} với lý thuyết. DP khớp hoàn toàn lý thuyết. Với AP, trừ xu thế làm giảm mạnh tự tương quan nhưng sai phân còn mùa vụ, nên ACF chưa đủ để kết luận. NI chỉ đổi mức, không thuộc cả hai loại.

\noidung{Bước ngẫu nhiên}
\tp{Nội dung lý thuyết cần kiểm chứng} Kết quả \ndlk{2.2.1(b)} nói rằng với bước ngẫu nhiên $Y_t=\delta+Y_{t-1}+\varepsilon_t$: $\Var Y_t=t\sigma^2$, $\Corr(Y_t,Y_{t+h})=\sqrt{t/(t+h)}$, ACF mẫu giảm chậm gần tuyến tính, và hồi quy giữa hai bước ngẫu nhiên độc lập cho kết quả giả mạo.
\tp{Cách triển khai} Ta kiểm tra gia số của cả ba chuỗi bằng Ljung--Box (10 độ trễ ở NI và DP, 12 ở AP). AP và DP dùng mẫu huấn luyện. Ta mô phỏng bước ngẫu nhiên có trôi với $(\hat\delta,\hat\sigma)$ của DP, với $B=2000$ đường dài $250$ ngày, để so phương sai, trung bình và tương quan với công thức; mô phỏng bước ngẫu nhiên không trôi $n=200$ cho ACF mẫu; hồi quy một bước ngẫu nhiên độc lập theo một bước khác với $n=100$ để đếm tần suất bác bỏ $H_0:\beta_1=0$ ở mức $5\%$. Cuối cùng là phần mở rộng đồng tích hợp: hồi quy log DAX theo log SMI và kiểm định Engle--Granger.
\tp{Mã lệnh Python}
\begin{codeblock}
d = np.diff(log_dax); delta, sigma = d.mean(), d.std(ddof=1)      # tham so tu DP
rng = S.child_rng(12, 2); B, T = 2000, 250
duong = np.array([S.simulate_random_walk(T, rng, delta, sigma) for _ in range(B)])
phuong_sai = duong[:, [0, 9, 49, 99, 249]].var(axis=0, ddof=1)    # so voi t * sigma^2
for _ in range(B):                                    # hoi quy gia mao, n = 100
    x, y = S.simulate_random_walk(n, rng), S.simulate_random_walk(n, rng)
    xc = x - x.mean(); b1 = (xc @ (y - y.mean())) / (xc @ xc)
    resid = (y - y.mean()) - b1 * xc
    se = np.sqrt(resid @ resid / (n - 2) / (xc @ xc))
    bac_bo += abs(b1 / se) > stats.t.ppf(0.975, n - 2)
\end{codeblock}
\tp{Kết quả tính toán} Bảng~\ref{tab:tab_b12_02_rw_variance} so phương sai mô phỏng với công thức, Bảng~\ref{tab:tab_b12_02_rw_correlation} và Bảng~\ref{tab:tab_b12_02_rw_sample_acf} cho tương quan và ACF mẫu của bước ngẫu nhiên, Bảng~\ref{tab:tab_b12_02_spurious_regression} cho tần suất bác bỏ của hồi quy giả mạo, và Bảng~\ref{tab:tab_b12_ext_cointegration} cho phần mở rộng đồng tích hợp.
\input{\tabdir/tab_b12_02_rw_variance.tex}
\input{\tabdir/tab_b12_02_rw_correlation.tex}
\input{\tabdir/tab_b12_02_rw_sample_acf.tex}
\input{\tabdir/tab_b12_02_spurious_regression.tex}
\input{\tabdir/tab_b12_ext_cointegration.tex}
\tp{Phân tích và diễn giải} Phương sai mô phỏng bám sát $t\hat\sigma^2$ ở mọi mốc (tỷ số $1.00$ đến $1.06$) và trung bình mô phỏng gần $\hat\delta t$. Tương quan mô phỏng $0.911$, $0.824$, $0.639$ gần $\sqrt{t/(t+h)}=0.913$, $0.816$, $0.632$ ở ba cặp $(t,h)$ được thử. ACF mẫu của bước ngẫu nhiên $n=200$ giảm chậm gần tuyến tính, với giá trị $0.963$, $0.927$, $0.825$, $0.675$ và $0.432$ ở các độ trễ 1, 2, 5, 10 và 20, nhỏ hơn $1$ vì $n$ hữu hạn. Hậu quả rõ nhất là hồi quy giả mạo: hai bước ngẫu nhiên độc lập mà kiểm định vẫn bác bỏ $\beta_1=0$ tới $76.9\%$ số lần, với sai số Monte Carlo $0.009$, gấp hơn $15$ lần mức $5\%$. Với phần mở rộng, hồi quy log DAX theo log SMI cho $\hat\beta=0.82$ và $R^2=0.965$, nhưng Engle--Granger cho $-2.50$ với $p=0.28$ không bác bỏ không đồng tích hợp: $R^2$ cao giữa hai chuỗi $I(1)$ không phải bằng chứng đồng tích hợp.

Nếu là bước ngẫu nhiên với gia số không tương quan, gia số phải vượt qua kiểm định Ljung--Box. NI cho $p=0.00078$ và AP cho $p<10^{-20}$, bác bỏ giả thiết này. DP cho $p=0.404$ trên mẫu huấn luyện, phù hợp về tự tương quan tuyến tính; phương sai thay đổi vẫn là hạn chế.

\tp{Nhận xét} Mô phỏng trên tham số DP \textbf{phù hợp} với lý thuyết: phương sai, tương quan, dạng ACF và hồi quy giả mạo đều khớp lý thuyết trong dung sai với tỷ số phương sai lệch tối đa $6.4\%$, khoảng $2$ lần sai số mô phỏng, và phần mở rộng cũng đúng với cảnh báo của lý thuyết.

\subsection{Mô hình ARIMA}

\noidung{Mô hình ARIMA và độ trôi}
\tp{Nội dung lý thuyết cần kiểm chứng} Kết quả \ndlk{2.2.2(a)} nói rằng ARIMA$(p,d,q)$ có $\nabla^dY_t$ là ARMA; hằng số khi $d=1$ là độ trôi $\mu_W=c/\phi(1)$, và bỏ độ trôi làm dự báo dài hạn nằm ngang thay vì đi theo xu thế.
\tp{Cách triển khai} Với AP, ta khớp ARIMA$(0,1,1)$ có trôi trên 120 tháng log hành khách và kiểm tra phần dư ở độ trễ 12. Ta khớp ARIMA$(0,1,1)$ không trôi cho NI, rồi ARIMA$(0,1,0)$ và ARIMA$(1,1,0)$ có trôi cho DP; so độ trôi và phương sai với giá trị tính tay từ chuỗi sai phân, so AIC, rồi dự báo 100 ngày của DP có và không có độ trôi.
\tp{Mã lệnh Python}
\begin{codeblock}
ni = ARIMA(NI, order=(0, 1, 1), trend="n").fit()               # Nile: khong troi
dp_rw = ARIMA(LD, order=(0, 1, 0), trend="t").fit()            # DAX: buoc ngau nhien + troi
dp_rw0 = ARIMA(LD, order=(0, 1, 0), trend="n").fit()           # khong co troi
dp_ar = ARIMA(LD, order=(1, 1, 0), trend="t").fit()
d1 = np.diff(LD); troi_tay, s2_tay = d1.mean(), d1.var(ddof=1) # so voi tinh tay
du_bao = [f.get_forecast(100).predicted_mean for f in (dp_rw, dp_rw0)]   # co / khong troi
\end{codeblock}
\tp{Kết quả tính toán} Bảng~\ref{tab:tab_b12_05_arima_params} cho tham số các mô hình đã khớp, Bảng~\ref{tab:tab_b12_05_rw_manual_vs_arima} so độ trôi và phương sai với giá trị tính tay, còn Hình~\ref{fig:fig_b12_05_drift_forecast} so dự báo có và không có độ trôi.
\input{\tabdir/tab_b12_05_arima_params.tex}
\input{\tabdir/tab_b12_05_rw_manual_vs_arima.tex}
\figx{fig_b12_05_drift_forecast}{DP: dự báo 100 ngày giao dịch của log giá DAX có và không có độ trôi}
\tp{Phân tích và diễn giải} Với NI, ARIMA$(0,1,1)$ cho $\hat\theta=-0.733$ với sai số chuẩn $0.074$ và $z=-9.9$, tức hằng số làm trơn mũ $\alpha=1+\hat\theta=0.267$ và AIC $1269.1$; mô hình không có trôi vì trung bình của chuỗi sai phân gần $0$, vì NI đổi mức chứ không có xu thế. Với DP, bước ngẫu nhiên có trôi cho độ trôi $6.47\times10^{-4}$ mỗi ngày với $z=2.66$ và $\hat\sigma^2=1.06\times10^{-4}$, rất gần giá trị tính tay là $6.52\times10^{-4}$ và $1.061\times10^{-4}$. Về lý thuyết, ước lượng hợp lý cực đại của độ trôi trong bước ngẫu nhiên bằng đúng trung bình của chuỗi sai phân; khoảng lệch $0.8\%$ ở đây đến từ tối ưu số trên một tham số rất nhỏ, chỉ bằng khoảng $0.02$ sai số chuẩn $2.4\times10^{-4}$, và ARIMA$(1,1,0)$ cho đúng $6.52\times10^{-4}$. Lưu ý quy ước của \texttt{statsmodels}: với \texttt{trend="t"} và $d=1$, hệ số \texttt{x1} là $\mu_W=\E\nabla Y_t$ chứ không phải hằng số $c=\phi(1)\mu_W$ trong \ndlk{2.2.2(a)}; hai đại lượng chỉ trùng nhau khi không có phần AR. Thêm hệ số AR không giúp gì: $\hat\phi=-0.0004$ với sai số chuẩn $0.018$, và AIC của bước ngẫu nhiên là $-11733$, thấp hơn $-11731$ của ARIMA$(1,1,0)$. Trên đồ thị, mô hình có trôi cho đường thẳng đi lên, cao hơn khoảng $0.065$ log-đơn vị sau $100$ ngày, còn mô hình bỏ trôi cho đường nằm ngang tại giá trị cuối, đúng như lý thuyết.

Ở AP, độ trôi ước lượng 0.00952 gần trung bình gia số 0.00926. Phần dư còn tự tương quan rất rõ ($p<10^{-18}$), nên mô hình diễn tả được độ trôi nhưng bỏ sót mùa vụ.

\tp{Nhận xét} Kết quả \textbf{phù hợp} với lý thuyết: độ trôi ước lượng rất gần tính tay, và việc có hay không có độ trôi làm dự báo dài hạn khác nhau đúng như lý thuyết. Độ trôi khác $0$ ở mức $5\%$ với $z=2.66$ nhưng chỉ đáng kể khi dự báo xa.

\noidung{Mô hình ARIMA mùa vụ}
\tp{Nội dung lý thuyết cần kiểm chứng} Kết quả \ndlk{2.2.2(b)} nói rằng nếu chuỗi $Y_t$ theo mô hình airline SARIMA$(0,1,1)(0,1,1)_{12}$ thì $\nabla\nabla_{12}Y_t$ là MA(13) có tự tương quan chỉ khác không ở các độ trễ $1$, $11$, $12$, $13$: $\rho(1)=\theta/(1+\theta^2)$, $\rho(12)=\Theta/(1+\Theta^2)$, $\rho(11)=\rho(13)=\rho(1)\rho(12)$.
\tp{Cách triển khai} Với NI ta thử chu kỳ 12 năm; với DP ta thử chu kỳ 5 phiên trên mẫu huấn luyện. Ta so ACF ở một và hai chu kỳ và tỷ số phương sai sau khi sai phân mùa vụ thêm một lần. Ta khớp airline cho log AP trên 120 tháng huấn luyện, so ACF mẫu của $\nabla\nabla_{12}$ log AP với ACF lý thuyết từ công thức đóng và từ khai triển MA(13); đồng thời so airline với 11 biến giả và hạng Fourier $K=2,3,4$ bằng AIC và RMSE trên phần kiểm định 1959--1960. Mỗi mô hình biến giả và Fourier được khớp hai lần: không có độ trôi, đúng như Ví dụ 4.2 của Bài 12, và có độ trôi, để tách tác động của xu thế khỏi cách mô hình hóa mùa vụ.
\tp{Mã lệnh Python}
\begin{codeblock}
d12 = TR[12:] - TR[:-12]; dd = np.diff(d12)                    # TR = log(AP) huan luyen
fit = airline(TR); par = dict(zip(fit.param_names, fit.params))
th, Th = par["ma.L1"], par["ma.S.L12"]
r_mau = A.sample_acf(dd, 25)                           # ACF mau cua chuoi da sai phan
ma = np.polymul([1, th], np.r_[1.0, np.zeros(11), Th])[1:]     # (1 + th z)(1 + Th z^12)
r_lt = R.theoretical_acf([], ma, 25, terms=60)          # ACF ly thuyet tu MA(13)
r1, r12 = th / (1 + th ** 2), Th / (1 + Th ** 2)
dong = {1: r1, 11: r1 * r12, 12: r12, 13: r1 * r12}     # cong thuc dong
\end{codeblock}
\tp{Kết quả tính toán} Bảng~\ref{tab:tab_b12_10_seasonal_lags} cho ACF ở độ trễ mùa vụ của ba chuỗi sai phân, Bảng~\ref{tab:tab_b12_10_airline_acf} và Hình~\ref{fig:fig_b12_10_airline_acf} so ACF mẫu với ACF lý thuyết của mô hình airline, còn Bảng~\ref{tab:tab_b12_10_seasonality_models} so airline với các cách mô hình hóa mùa vụ khác.
\input{\tabdir/tab_b12_10_seasonal_lags.tex}
\input{\tabdir/tab_b12_10_airline_acf.tex}
\figx{fig_b12_10_airline_acf}{ACF của $\nabla\nabla_{12}$ log AP (que) và ACF lý thuyết của mô hình airline (chấm)}
\input{\tabdir/tab_b12_10_seasonality_models.tex}
\tp{Phân tích và diễn giải} Độ lệch chuẩn của phần huấn luyện giảm đều qua từng phép sai phân, từ $0.392$ ở log AP xuống $0.106$ sau $\nabla$, $0.066$ sau $\nabla_{12}$ và $0.046$ sau $\nabla\nabla_{12}$, không có dấu hiệu sai phân thừa. ACF mẫu khớp ACF lý thuyết ở các độ trễ đặc trưng: $-0.328$ so với $-0.307$ ở độ trễ 1, $-0.406$ so với $-0.418$ ở độ trễ 12; ở độ trễ 11 và 13 là $0.076$ và $0.182$ so với $0.128$, trong phạm vi sai số mẫu khoảng $0.1$. Công thức đóng và khai triển MA(13) cho cùng giá trị. Ở các độ trễ còn lại lý thuyết cho $0$ và phần lớn ACF mẫu nằm trong dải $\pm0.189$; chỉ độ trễ 3 với $-0.22$ và độ trễ 23 với $0.23$ vượt dải, hợp lý khi xét $25$ độ trễ. Về dự báo, khi khớp đúng như Ví dụ 4.2, tức không có độ trôi, airline cho RMSE $43.2$ trên phần kiểm định, tốt hơn rõ so với $85.5$ của 11 biến giả và $74.8$ đến $82.0$ của Fourier. Nhưng các mô hình này chỉ lấy một lần sai phân và không có độ trôi, nên dự báo dài hạn nằm ngang, trong khi airline với $\nabla\nabla_{12}$ tự mang xu thế tuyến tính. Khi thêm độ trôi, được ước lượng khoảng $0.0093$ đến $0.0099$ mỗi tháng, sát trung bình $0.0093$ của $\nabla\log$ AP, RMSE giảm còn $33.1$ với 11 biến giả và từ $21.9$ đến $27.6$ với Fourier, tức tốt hơn airline. Như vậy khoảng cách ở Ví dụ 4.2 chủ yếu đến từ việc có hay không có xu thế chứ không phải từ cách mô hình hóa mùa vụ; trên 24 tháng kiểm định, mô hình mùa vụ cố định có độ trôi còn dự báo tốt hơn airline. AIC không so sánh trực tiếp được giữa các mô hình có cách sai phân khác nhau, nên AIC thấp hơn của mô hình biến giả, $-424.3$ hay $-434.0$ so với $-389.0$, không cho biết mô hình nào tốt hơn.

Ở NI, ACF tại 12 và 24 năm là $-0.044,-0.040$ (dải $\pm0.197$). Ở DP, ACF tại 5 và 10 phiên là $-0.011,-0.014$ (dải $\pm0.051$). Phép sai phân theo chu kỳ thử làm phương sai tăng 2.03 và 2.02 lần. Hai chuỗi chỉ xác nhận bước chẩn đoán; dữ liệu chưa ủng hộ việc dùng sai phân mùa vụ ở các chu kỳ này.

\tp{Nhận xét} Riêng AP \textbf{phù hợp} với mô hình mùa vụ: cấu trúc bốn độ trễ của ACF được xác nhận cả về giá trị lẫn vị trí, với một mô hình chỉ có hai tham số; hai độ trễ 3 và 23 vượt dải là ngoại lệ nhỏ. Về dự báo, airline vượt các mô hình mùa vụ không có độ trôi như trong Ví dụ 4.2, nhưng không vượt được khi các mô hình đó có độ trôi. Trên 24 tháng kiểm định, thứ hạng dự báo vì vậy phụ thuộc vào cách mô hình hóa xu thế.

\subsection{Xây dựng mô hình ARIMA}

\noidung{Sai phân và sai phân quá mức}
\tp{Nội dung lý thuyết cần kiểm chứng} Kết quả \ndlk{2.2.3(a)} nói rằng sai phân một chuỗi dừng tạo ra sai phân quá mức: MA(1) với $\theta=-1$, $\rho(1)=-\tfrac12$, không khả nghịch, phương sai dài hạn bằng $0$; sai phân thừa làm dự báo kém đi gấp đôi, sai phân thiếu có thể tệ hơn nhiều.
\tp{Cách triển khai} Với AP ta so chuỗi log sau sai phân thường và mùa vụ với chính chuỗi đó sau khi sai phân thường thêm một lần. Trên dữ liệu thực, ta so sai phân một và hai lần qua $\hat\rho(1)$, phương sai dài hạn ước lượng theo Bartlett và $\hat\theta$ của MA(1) khớp cho chuỗi đã sai phân. Bằng mô phỏng, ta kiểm tra $\hat\rho(1)$ của nhiễu trắng đã sai phân với $1000$ chuỗi độ dài $n=400$, và hai chi phí với $B=5000$ lần lặp và $n=100$: dự báo nhiễu trắng bằng bước ngẫu nhiên, dự báo bước ngẫu nhiên bằng trung bình mẫu.
\tp{Mã lệnh Python}
\begin{codeblock}
for ten, y in (("NI", NI), ("DP", LD)):
    for nhan, w, xu_the in (("d1", np.diff(y), "c"), ("d2 (qua muc)", np.diff(y, 2), "n")):
        m = ARIMA(w, order=(0, 0, 1), trend=xu_the).fit()
        rho1, theta, pv_dai_han = A.sample_acf(w, 1)[1], m.maparams[0], long_run_var(w)
    for K in (3, 7, 15, 31, 63):                                # Bartlett theo K, so voi gia tri dong
        g = A.sample_autocov(np.diff(y), K + 1)
        bartlett_K, gia_tri_dong = long_run_var(np.diff(y, 2), K), 2 * (g[0] - g[K + 1]) / (K + 1)
e = rng.normal(size=n + 1); e_rw.append(e[n] - e[n - 1])      # nhieu trang, du bao bang Y_n
w = np.cumsum(rng.normal(size=n + 1))                         # buoc ngau nhien
e_tb.append(w[n] - w[:n].mean())                              # du bao bang trung binh mau
\end{codeblock}
\tp{Kết quả tính toán} Bảng~\ref{tab:tab_b12_03_overdifferencing} cho kết quả sai phân một và hai lần trên dữ liệu thực, Bảng~\ref{tab:tab_b12_03_bartlett_by_k} cho ước lượng Bartlett của phương sai dài hạn theo số độ trễ $K$, còn Bảng~\ref{tab:tab_b12_03_cost_simulation} cho chi phí của sai phân thừa và thiếu qua mô phỏng.
\input{\tabdir/tab_b12_03_overdifferencing.tex}
\input{\tabdir/tab_b12_03_bartlett_by_k.tex}
\input{\tabdir/tab_b12_03_cost_simulation.tex}
\tp{Phân tích và diễn giải} Với NI, sai phân một lần cho $\hat\rho(1)=-0.40$ và $\hat\theta=-0.765$, phù hợp với ARIMA$(0,1,1)$ ở \muc{5.2}; giá trị này khác $-0.733$ ở \ndlk{5.2(a)} vì ở đây MA(1) được khớp cho $\nabla$NI có hằng số, tức có độ trôi, còn ở \ndlk{5.2(a)} mô hình không có độ trôi; sai phân lần hai đưa $\hat\rho(1)$ xuống $-0.63$, $\hat\theta=-0.999$ sát biên với sai số chuẩn $2.1$, và độ lệch chuẩn tăng từ $168$ lên $283$: đây là sai phân thừa rõ ràng. Với DP, sai phân một lần cho gần nhiễu trắng ($\hat\rho(1)=-0.0004$); sai phân lần hai tạo đúng sai phân thừa của lý thuyết: $\hat\rho(1)=-0.486$ sát $-\tfrac12$, $\hat\theta=-0.9986$, độ lệch chuẩn tăng từ $0.0103$ lên $0.0146$, xấp xỉ $\sqrt2$ lần. Phương sai dài hạn ước lượng theo Bartlett giảm ở DP, từ $9.7\times10^{-5}$ xuống $2.7\times10^{-5}$, nhưng chưa về $0$, còn ở NI thì tăng từ $10252$ lên $15391$. Điều này đến từ chính ước lượng Bartlett chứ không mâu thuẫn với lý thuyết. Với $W=\nabla X$, tổng Bartlett $K$ độ trễ $\sum_{|h|\le K}\bigl(1-\tfrac{|h|}{K+1}\bigr)\gamma_W(h)$ bằng đúng $2\bigl(\gamma_X(0)-\gamma_X(K+1)\bigr)/(K+1)$, vì $\gamma_W(h)=2\gamma_X(h)-\gamma_X(h-1)-\gamma_X(h+1)$ và sai phân bậc hai của trọng số tam giác chỉ khác $0$ tại $h=0$ và $h=\pm(K+1)$. Đại lượng này giảm như $1/(K+1)$ chứ không bằng $0$ ở $K$ hữu hạn. Khi $X$ là nhiễu trắng, như $\nabla\log$ DAX, nó bằng $2\sigma^2/(K+1)=\sigma^2/4$ với $K=7$, khớp với mức $2.7\times10^{-5}\approx\hat\sigma^2/4$. Thay $\gamma_X$ bằng ước lượng mẫu, công thức cho $15221$ ở NI và $2.67\times10^{-5}$ ở DP, sát hai giá trị trong bảng. Bảng~\ref{tab:tab_b12_03_bartlett_by_k} cho thấy khi tăng $K$, ước lượng của chuỗi sai phân thừa giảm đúng theo $1/(K+1)$ về $0$. Riêng NI, $\hat\rho(1)=-0.63$ của $\nabla^2$NI nằm ngoài $[-\tfrac12,\tfrac12]$, miền cho phép của MA(1) theo \ndlk{2.1.2(d)}. Lý do là $\nabla$NI vốn đã có $\hat\rho(1)=-0.40$, nên $\nabla^2$NI gần một MA(2) có nghiệm đơn vị chứ không phải MA(1); $\hat\theta=-0.999$ với sai số chuẩn $2.1$ vì vậy phản ánh cả việc khớp sai mô hình lẫn sai phân thừa. Mô phỏng xác nhận phần còn lại, đều trong dung sai $3$ MCSE: trung bình $\hat\rho(1)$ của nhiễu trắng đã sai phân là $-0.497$, so với lý thuyết là $-0.5$, lệch $2.6$ MCSE; dự báo nhiễu trắng bằng bước ngẫu nhiên có phương sai sai số $2.005$, so với lý thuyết là $2$ và gấp đôi mức tối ưu $1$, lệch $0.1$ MCSE; dự báo bước ngẫu nhiên bằng trung bình mẫu có phương sai $33.7$, so với lý thuyết là $33.8$ và gấp $33$ lần mức tối ưu, lệch $0.2$ MCSE.

Trên AP, sai phân thường bổ sung làm ACF bậc một giảm từ $-0.328$ xuống $-0.660$. MA(1) trên chuỗi mới cho $\hat\theta=-0.99995$, môđun nghiệm 1.00005. Nghiệm gần biên là dấu hiệu sai phân thừa; giá trị $-1/2$ chỉ là mốc cho trường hợp chuỗi trước sai phân là nhiễu trắng.

\tp{Nhận xét} Kết quả \textbf{phù hợp} với lý thuyết. Dấu hiệu chính của sai phân thừa ($\hat\rho(1)\to-0.5$, $\hat\theta\to-1$) và hai chi phí được xác nhận chính xác. Phương sai dài hạn ước lượng chưa về $0$ ở $K=3$ (NI) và $K=7$ (DP) đúng như giá trị $2(\gamma_X(0)-\gamma_X(K+1))/(K+1)$ của ước lượng Bartlett, và giảm về $0$ theo $1/(K+1)$ khi tăng $K$, nên đây không phải điểm lệch so với lý thuyết. Với NI cần lưu ý thêm: $\hat\rho(1)$ vượt $-\tfrac12$ cho thấy $\nabla^2$NI không phải MA(1) đơn thuần, nên ngoài sai phân thừa còn có sai mô hình.

\noidung{Kiểm định nghiệm đơn vị}
\tp{Nội dung lý thuyết cần kiểm chứng} Kết quả \ndlk{2.2.3(b)} nói rằng thống kê Dickey--Fuller có phân phối giới hạn không chuẩn, dịch về phía âm, với giá trị tới hạn 5\% là $-2.86$ khi có hệ số chặn; ADF với giả thuyết không là nghiệm đơn vị và KPSS với giả thuyết không là dừng nên đọc cùng nhau; công suất của ADF thấp khi $\phi$ gần $1$ hoặc $n$ nhỏ.
\tp{Cách triển khai} Nội dung này dùng cả dữ liệu thực lẫn mô phỏng, và áp dụng cho NI, AP và DP. Trên dữ liệu thực, ta chạy ADF với độ trễ chọn theo AIC và KPSS cho chuỗi gốc, có hệ số chặn và có cả xu thế, và cho sai phân bậc một. Bằng mô phỏng, ta lấy phân phối của $\tau$ dưới $H_0$ từ $B=2000$ bước ngẫu nhiên với $n=200$, so với $-2.86$ và $-1.645$, rồi đo tỷ lệ bác bỏ với $B=1000$ lần mô phỏng AR(1) có $\phi=0.95,\,0.9,\,0.8$ ở $n=100$ và $200$, và với nhiễu iid cộng dịch mức. Tỷ lệ bác bỏ được tính trên cùng các chuỗi mô phỏng cho hai cấu hình: DF không độ trễ, và ADF chọn độ trễ theo AIC như mã mẫu của bài giảng.
\tp{Mã lệnh Python}
\begin{codeblock}
def doc(y, kiem_dinh, hoi_quy):
    if kiem_dinh == "ADF":
        s, p, lag, _, crit, _ = adfuller(y, regression=hoi_quy, autolag="AIC")
        return s, p, crit["5%"], lag
    s, p, lag, crit = kpss(y, regression=hoi_quy, nlags="auto")
    return s, p, crit["5%"], lag
tau = np.array([adfuller(S.simulate_random_walk(200, rng), maxlag=0,       # tau duoi H0
                         regression="c", autolag=None)[0] for _ in range(2000)])
ty_le_dung, ty_le_sai = np.mean(tau < -2.86), np.mean(tau < -1.645)       # 0.05 / bang chuan
\end{codeblock}
\tp{Kết quả tính toán} Bảng~\ref{tab:tab_b12_04_unit_root_tests} cho kết quả ADF và KPSS trên dữ liệu thực, Bảng~\ref{tab:tab_b12_04_null_distribution} cho phân phối của $\tau$ dưới giả thuyết không, và Bảng~\ref{tab:tab_b12_04_adf_power} cho công suất của ADF.
\input{\tabdir/tab_b12_04_unit_root_tests.tex}
\input{\tabdir/tab_b12_04_null_distribution.tex}
\input{\tabdir/tab_b12_04_adf_power.tex}
\tp{Phân tích và diễn giải}
\bodl{NI.} ADF bác bỏ nghiệm đơn vị với $-4.05$ và $p=0.001$, \emph{và} KPSS cũng bác bỏ tính dừng với $0.869>0.463$: đây là khuôn mẫu ``cùng bác bỏ'' mà lý thuyết gắn với điểm gãy, ở đây là lần đổi mức quanh 1899. Sau sai phân, ADF bác bỏ với $-4.68$ và KPSS không bác bỏ với $0.117$. Hai kết quả ủng hộ tính dừng của chuỗi sai phân.

\bodl{AP.} ADF không bác bỏ với $p=0.32$, và $p=0.80$ khi có xu thế, còn KPSS bác bỏ với $1.744>0.463$, nên chuỗi nghiêng về $I(1)$; nhưng KPSS có xu thế không bác bỏ với $0.061<0.146$, nên chưa loại trừ được mức dừng quanh xu thế. Sau sai phân bậc một, ADF vẫn không bác bỏ với $p=0.158$, nhiều khả năng do mùa vụ chưa được loại bỏ, vì ACF ở độ trễ 12 vẫn lớn như đã thấy ở \mucnd{5.1(a)}.

\bodl{DP.} ADF không bác bỏ với $p=0.996$ và KPSS bác bỏ rõ với $5.75$, nên hai kiểm định thống nhất ủng hộ nghiệm đơn vị. Sau sai phân, ADF bác bỏ mạnh với $-43.1$ và KPSS không bác bỏ với $0.437<0.463$, nên sai phân bậc một là đủ.

\bodl{Mô phỏng.} Dưới $H_0$ với $n=200$, phân vị $5\%$ của $\tau$ là $-2.79$ so với $-2.86$ trong bảng, trung bình là $-1.51$ so với $-1.53$, và tần suất $\tau<-2.86$ là $4.2\%$ so với $5\%$ lý thuyết, lệch $1.7$ lần sai số Monte Carlo. Nếu dùng giá trị tới hạn chuẩn $-1.645$, tần suất bác bỏ sai lên $45.7\%$, hơn chín lần mức danh nghĩa. Công suất thấp khi $\phi$ gần $1$: với DF không độ trễ, $\phi=0.95$ chỉ cho $13.8\%$ ở $n=100$ và $32.7\%$ ở $n=200$; $\phi=0.9$ cho $33.1\%$ và $88.6\%$; $\phi=0.8$ cho $89.6\%$ và $100\%$. ADF chọn độ trễ theo AIC cho các mức tương tự, từ $16.7\%$ đến $98.3\%$, sát các giá trị $0.14$ đến $0.99$ của bài giảng. Khác biệt lớn nằm ở nhiễu iid cộng dịch mức: DF không độ trễ bác bỏ $100\%$ số lần, còn ADF chọn độ trễ theo AIC chỉ bác bỏ $8\%$, đúng bằng con số $0.08$ của bài giảng. Đây là hiện tượng Perron: bước dịch mức tạo tự tương quan dương kéo dài trong chuỗi, AIC bù bằng cách thêm độ trễ, và kiểm định gần như mất hết công suất dù chuỗi không có nghiệm đơn vị. Khi bắt buộc không có độ trễ, hồi quy DF chỉ thấy $\hat\rho(1)\approx0.5$, xa $1$, nên bác bỏ. Tỷ lệ $8\%$ chỉ mô tả thiết kế mô phỏng trên. Với chuỗi NI thực tế, AIC chọn một độ trễ và ADF bác bỏ nghiệm đơn vị ($p=0.001$) dù chuỗi có đổi mức. Vì vậy khi nghi có dịch mức, kết quả ADF cần được đọc cùng KPSS và đồ thị của chuỗi, như đã làm với NI ở trên.

\tp{Nhận xét} Kết quả \textbf{phù hợp} với lý thuyết: phân phối của $\tau$ khớp bảng Dickey--Fuller và khác hẳn phân phối chuẩn; công suất thấp gần nghiệm đơn vị; trong mô phỏng có dịch mức, nhiễu dừng quanh từng mức trung bình và ADF chọn độ trễ theo AIC ít bác bỏ nghiệm đơn vị; và việc đọc ADF cùng KPSS giải thích đúng từng chuỗi, nổi bật là Nile. Với AP kết quả chưa dứt khoát vì mùa vụ và xu thế đan xen, cần kiểm định riêng cho mùa vụ.
